\frontmatterpage
\unchapter{Abstract}

This Bachelor's thesis examined the integration maturity of information systems
within Project Management Office (PMO) environments in medium and large Latvian
enterprises. As organizations increasingly rely on multiple digital tools for
project planning, reporting, and governance, the degree to which these systems
are integrated has a direct impact on data quality, reporting efficiency, and
governance standardization. Despite its practical relevance, this topic had
received limited empirical attention in the Latvian context.

The Bachelor Paper consists of five chapters. Chapter 1 reviewed the theoretical
background, defining key concepts such as PMO, information system integration,
maturity models, business intelligence, and data quality, and surveyed existing
research on PMO digitalization and integration maturity frameworks. Chapter 2
presented the research methodology, including the design of a four-level
Integration Maturity Model and the development of a structured survey instrument.
Chapter 3 contained the empirical study, presenting the results of a survey
conducted among representatives of Latvian enterprises with formal PMO structures,
supplemented by one structured expert interview and statistical analysis. Chapter 4
discussed the findings in the context of existing literature and proposed five
practical recommendations for improving information system integration maturity
in PMO environments. Chapter 5 reflected on future research directions and the
broader societal and economic implications of the study.

The methodological basis of the study included systematic literature review,
survey-based quantitative research, one structured expert interview, descriptive
statistical analysis, and Spearman rank correlation analysis.

The study found that 70\% of PMO environments in Latvia operate at intermediate
integration maturity levels (Levels 2 or 3), and that higher integration maturity
is significantly associated with greater reporting automation (r=0.405, p=0.010),
better data quality (r=0.321, p=0.043), and stronger governance standardization
(r=0.403, p=0.010). All four research hypotheses were supported. The findings
provide evidence-based guidance for organizations seeking to improve their project
management information systems.

\thesisscopeEN

\textbf{Keywords:} Project Management Office, Information System Integration,
Integration Maturity Model, Reporting Automation, Data Quality, Business
Intelligence, Governance Standardization.
